    \section{Моделирование потери GNSS на наборе EuRoC}\label{sec:gnss_loss}

    Стандартные метрики главы~\ref{sec:euroc} характеризуют точность всей траектории, совмещённой с эталоном задним числом.
    Для задачи работы важнее другой вопрос: насколько быстро растёт ошибка положения, если до некоторого момента навигация опиралась на GNSS, а после него -- только на визуально-инерциальную оценку.
    В этой главе описана методика, отвечающая на этот вопрос, и приведены результаты её применения к запускам главы~\ref{sec:euroc}.

    \subsection{Методика}

    \textbf{Сценарий.}
    Пока GNSS работает, автопилот знает положение и курс аппарата и может определить, как система координат визуально-инерциальной оценки связана с глобальной.
    В момент $t_0$ сигнал GNSS пропадает, и дальше положение вычисляется только по визуально-инерциальной оценке с этой последней известной связью.
    Ошибка навигации после $t_0$ -- это расстояние между положением, полученным таким образом, и истинным положением.

    Рассматриваемые системы не используют GNSS, поэтому его потеря не меняет их работу: меняется только то, как их оценка привязывается к глобальной системе координат.
    Это позволяет не перезапускать системы для каждого момента отказа, а моделировать отказ при оценке: из одного запуска получается множество экспериментов с разными $t_0$.
    Отказ моделируется в моменты $t_0$ с шагом $5$~с по всему полёту, а результаты усредняются по всем $t_0$ и всем пяти запускам.

    \textbf{Онлайн-оценка.}
    Итоговые траектории ORB-SLAM3 и траектория VINS-Fusion после замыкания петель сохраняются в конце работы, после глобальной оптимизации, поэтому каждая поза в них уточнена с учётом последующих кадров.
    На борту таких уточнений в момент полёта нет, поэтому для моделирования потери GNSS использовалась \emph{онлайн-оценка} -- поза, которую система выдала в момент обработки кадра.
    Для ORB-SLAM3 в пример запуска добавлена запись позы, возвращаемой при обработке каждого кадра, вместе с состоянием трекинга; для VINS-Fusion записывается поза VIO, исправленная текущим графом поз (топик \texttt{odometry\_rect}).
    Выход VIO VINS-Fusion и OpenVINS и так формируется по ходу работы и используется без изменений.

    \textbf{Привязка к эталону до отказа.}
    Связь систем координат определяется по окну $[t_0 - W,\, t_0]$, $W = 10$~с.
    Совмещение только по положениям (как в главе~\ref{sec:euroc}) плохо обусловлено, если аппарат в окне зависает или летит прямо, поэтому поворот находится по ориентациям:
    \begin{equation}\label{eq:window_rot}
        R = \arg\min_{R \in SO(3)} \sum_{i \in \mathcal{W}} \left\| R_i^{\mathrm{ref}} - R\, R_i^{\mathrm{est}} \right\|_F^2,
        \qquad
        d = \frac{1}{|\mathcal{W}|} \sum_{i \in \mathcal{W}} \left( p_i^{\mathrm{ref}} - R\, p_i^{\mathrm{est}} \right),
    \end{equation}
    где $\mathcal{W}$ -- позы в окне, $R_i, p_i$ -- ориентации и положения эталона и оценки, $R, d$ -- искомые поворот и сдвиг.
    Задача для $R$ решается через сингулярное разложение матрицы $\sum_i R_i^{\mathrm{ref}} (R_i^{\mathrm{est}})^\top$.
    Ошибка в момент $t > t_0$ равна $e(t) = \left\| p^{\mathrm{ref}}(t) - (R\, p^{\mathrm{est}}(t) + d) \right\|$ и рассматривается как функция времени $t - t_0$ и пути, пройденного после $t_0$.

    \textbf{Разрывы оценки.}
    Онлайн-оценка может прерываться: система перестаёт выдавать позу (потеря трекинга, сброс карты) или скачком переводит карту в другую систему координат (завершение инерциальной инициализации, уточнение масштаба, слияние карт).
    Траектория делится на непрерывные участки; граница участка -- пауза в выдаче поз дольше $0{,}25$~с или изменение положения со скоростью больше $5$~м/с (скорость аппарата в EuRoC не превышает $2{,}5$~м/с).
    Окно привязки и интервал после $t_0$ берутся с одного участка: если внутри окна был разрыв, момент $t_0$ не рассматривается, а после разрыва навигация по этой оценке считается потерянной.
    Доля полёта, покрытая самым длинным непрерывным участком, приводится отдельно.

    \textbf{Метрики.}
    Для каждой системы и последовательности вычисляется медиана по всем $t_0$ ошибки через $10$ и $30$~с и через $10$ и $25$~м пути после потери GNSS.
    Медиана выбрана вместо среднего, так как отдельные $t_0$, приходящиеся на резкие манёвры, дают редкие большие ошибки.
    Более длинные горизонты на EuRoC достигаются лишь из небольшого числа $t_0$, особенно для ORB-SLAM3, поэтому в таблицы не включены.

    \subsection{Результаты}

    Зависимость ошибки от времени после потери GNSS показана на рис.~\ref{fig:gnss_loss}, значения на фиксированных горизонтах -- в таблицах~\ref{table:gnss_loss_time} и~\ref{table:gnss_loss_dist}, доля полёта с непрерывной онлайн-оценкой -- в таблице~\ref{table:online_coverage}.

    \begin{figure}[H]
        \centering
        \includegraphics[width=\textwidth]{euroc_gnss_loss.png}
        \caption{Ошибка положения после потери GNSS (медиана по моментам отказа и пяти запускам)}
        \label{fig:gnss_loss}
    \end{figure}

    \begin{table}[H]
        \caption{Ошибка положения через $10$ и $30$~с после потери GNSS, см (медиана по моментам отказа)}
        \label{table:gnss_loss_time}
        \small
        \centering
        \renewcommand{\arraystretch}{1.2}
        \begin{tabular}{|l|c|c|c|c|c|}
            \hline
            \textbf{Метод} & \textbf{MH01} & \textbf{MH03} & \textbf{MH05} & \textbf{V102} & \textbf{V202} \\
            \hline
            ORB-SLAM3 & 7 / 8 & 10 / 8 & 12 / 15 & 9 / 7 & 3 / 2 \\
            VINS-Fusion (VIO) & 10 / 15 & 18 / 22 & 31 / 48 & 13 / 19 & 11 / 14 \\
            VINS-Fusion с замыканием петель & 11 / 17 & 18 / 16 & 34 / 51 & 16 / 15 & 10 / 12 \\
            OpenVINS & 8 / 12 & 13 / 17 & 20 / 72 & 8 / 13 & 6 / 8 \\
            \hline
        \end{tabular}
    \end{table}

    \begin{table}[H]
        \caption{Ошибка положения через $10$ и $25$~м пути после потери GNSS, см (медиана по моментам отказа)}
        \label{table:gnss_loss_dist}
        \small
        \centering
        \renewcommand{\arraystretch}{1.2}
        \begin{tabular}{|l|c|c|c|c|c|}
            \hline
            \textbf{Метод} & \textbf{MH01} & \textbf{MH03} & \textbf{MH05} & \textbf{V102} & \textbf{V202} \\
            \hline
            ORB-SLAM3 & 9 / 9 & 9 / 11 & 12 / 16 & 4 / 7 & 3 / 2 \\
            VINS-Fusion (VIO) & 15 / 25 & 22 / 25 & 31 / 45 & 11 / 18 & 10 / 13 \\
            VINS-Fusion с замыканием петель & 19 / 30 & 19 / 21 & 36 / 46 & 19 / 16 & 10 / 11 \\
            OpenVINS & 9 / 14 & 12 / 16 & 25 / 55 & 9 / 9 & 6 / 9 \\
            \hline
        \end{tabular}
    \end{table}

    \begin{table}[H]
        \caption{Доля полёта, покрытая самым длинным непрерывным участком онлайн-оценки, \% (среднее по пяти запускам)}
        \label{table:online_coverage}
        \small
        \centering
        \renewcommand{\arraystretch}{1.2}
        \begin{tabular}{|l|c|c|c|c|c|}
            \hline
            \textbf{Метод} & \textbf{MH01} & \textbf{MH03} & \textbf{MH05} & \textbf{V102} & \textbf{V202} \\
            \hline
            ORB-SLAM3 & 65 & 63 & 57 & 74 & 70 \\
            VINS-Fusion (VIO) & 100 & 99 & 97 & 95 & 97 \\
            VINS-Fusion с замыканием петель & 100 & 80 & 87 & 95 & 97 \\
            OpenVINS & 76 & 87 & 82 & 95 & 98 \\
            \hline
        \end{tabular}
    \end{table}

    \subsection{Анализ результатов}

    \textbf{ORB-SLAM3.}
    Ошибка ORB-SLAM3 после потери GNSS наименьшая на всех последовательностях и почти не растёт со временем: через $30$~с она составляет $2$--$15$~см.
    Причина в характере данных: аппарат в EuRoC многократно пролетает над одними и теми же участками небольшого помещения, и система локализуется относительно карты, построенной ещё до отказа, а не накапливает ошибку одометрии.
    При полёте над новой местностью, типичном для поисковой миссии, этого преимущества не будет, поэтому результат ORB-SLAM3 на EuRoC нельзя переносить на полёт БПЛА без проверки на соответствующих данных.

    Обратная сторона -- долгий выход на устойчивую оценку.
    Непрерывная онлайн-оценка ORB-SLAM3 покрывает только $57$--$74$\% полёта: в начале система многократно сбрасывает карту, пока не выполнит инерциальную инициализацию, а затем при уточнении масштаба поза меняется скачками (в одном из предварительных запусков на MH01 -- на $64$~м).
    В итоговой траектории этих скачков не видно, так как она пересчитана после оптимизации.
    Если GNSS пропадёт раньше, чем завершится инициализация, ORB-SLAM3 не даст пригодной для навигации оценки.

    \textbf{Одометрия без карты.}
    OpenVINS и VIO-часть VINS-Fusion накапливают ошибку со временем, и именно они характеризуют поведение при полёте над новой местностью.
    OpenVINS точнее: через $30$~с его ошибка составляет $8$--$17$~см на четырёх последовательностях против $14$--$22$~см у VINS-Fusion.
    На MH05 (быстрые манёвры и тёмные участки) ошибка OpenVINS через $30$~с вырастает до $72$~см, у VINS-Fusion -- до $48$~см.
    У VINS-Fusion ошибка растёт примерно линейно: на MH01 -- $10$, $15$ и $29$~см через $10$, $30$ и $60$~с.

    \textbf{Замыкание петель.}
    Онлайн-оценка VINS-Fusion с замыканием петель после потери GNSS в целом не точнее чистого VIO: на коротких горизонтах ошибка близка или больше, на длинных на части последовательностей меньше (рис.~\ref{fig:gnss_loss}).
    Исправление после замыкания петли приходит с задержкой и сдвигает оценку к системе координат начала полёта, тогда как привязка по окну до $t_0$ уже учла накопленный к этому моменту дрейф.
    Кроме того, крупные исправления воспринимаются как скачки: на MH03 в одном из запусков они сократили непрерывный участок до половины полёта.
    Таким образом, уменьшение ошибки итоговой траектории в $1{,}3$--$2{,}9$ раза за счёт петель (глава~\ref{sec:euroc}) в момент полёта не реализуется, и для оценки систем по критерию задачи работы необходимо использовать онлайн-оценку.

    \subsection{Ограничения}

    \begin{itemize}
        \item Набор EuRoC снят в помещении, полёты длятся $85$--$184$~с со скоростью до $2{,}5$~м/с на высоте нескольких метров; повторные пролёты над одними участками благоприятствуют системам с картой.
        \item При привязке по окну до $t_0$ используется точный эталон; реальный приёмник GNSS имеет погрешность порядка метров, и ошибка привязки, особенно по курсу, увеличит ошибку после отказа. Полученные значения -- нижняя оценка.
        \item Моменты $t_0$ одного запуска с близкими временами дают коррелированные результаты, поэтому медиана по ним точнее отражает типичную ошибку, чем её разброс.
    \end{itemize}

    \subsection{Выводы по главе}

    Разработана методика оценки накопления ошибки после потери GNSS, использующая онлайн-оценку систем, привязку к глобальной системе координат по окну до отказа и учёт разрывов оценки.
    Применение методики к $75$ запускам на наборе EuRoC показало, что итоговые траектории и стандартные метрики дают завышенное представление о пригодности систем для навигации: ORB-SLAM3 выдаёт устойчивую онлайн-оценку лишь на $57$--$74$\% полёта, а замыкание петель не уменьшает ошибку VINS-Fusion в момент полёта.
    Среди систем без опоры на ранее построенную карту наименьшую ошибку после потери GNSS показал OpenVINS (через $30$~с -- $8$--$17$~см, кроме MH05).
    Для оценки поведения систем при полёте над новой местностью на высоте и скорости, характерных для БПЛА, требуются данные, приближенные к этим условиям.
    \clearpage
